\chapter{Non-Sinusoidal Signals and Fourier Analysis}
\label{ch:non-sine}

Not every useful waveform is a sine wave. Square waves clock digital
logic, sawtooth waves sweep oscilloscopes and older TV displays, and the
keyed on/off waveform of Morse code is about as far from a sine wave as a
signal can get. This chapter explains how such waveforms relate back to
the sinusoidal theory of Chapter~\ref{ch:ac-theory}, and why that
relationship matters for bandwidth and interference.

\section{Common Non-Sinusoidal Waveforms}

\begin{figure}[H]
\centering
\begin{tikzpicture}
\begin{groupplot}[
  group style={group size=2 by 2, horizontal sep=1.6cm, vertical sep=1.4cm},
  width=6.2cm, height=3.6cm, axis lines=middle,
  xtick=\empty, ytick=\empty, ymin=-1.4, ymax=1.4, xmin=0, xmax=4,
  domain=0:4, samples=200,
  every axis plot/.append style={line width=1pt},
  every axis x label/.style={at={(axis description cs:0.5,-0.05)},anchor=north},
]
\nextgroupplot[title={Sine},xlabel={$t$}]
\addplot[blue] {sin(x*180)};

\nextgroupplot[title={Square},xlabel={$t$}]
\addplot[red,samples=400,unbounded coords=jump] {sign(sin(x*180))};

\nextgroupplot[title={Triangle},xlabel={$t$}]
\addplot[rorange] {asin(sin(x*180))/90};

\nextgroupplot[title={Sawtooth},xlabel={$t$}]
\addplot[rgreenline,samples=400,unbounded coords=jump] {2*(x/2-floor(x/2+0.5))};
\end{groupplot}
\end{tikzpicture}
\caption{Common periodic waveforms: sine, square, triangle, and
sawtooth, each shown over several cycles.}
\label{fig:waveforms}
\end{figure}

\begin{itemize}
\item \textbf{Square wave}: switches abruptly between two levels, spending
equal time at each (50\% duty cycle). Produced by digital logic gates,
keyed CW (approximately), and switching-mode power supplies.
\item \textbf{Triangle wave}: rises and falls linearly at a constant
rate. Used in some function generators and as a building block for
sweep and ramp circuits.
\item \textbf{Sawtooth wave}: rises linearly then drops sharply (or vice
versa). Classic use: the horizontal deflection sweep in a
cathode-ray-tube display.
\item \textbf{Pulse train}: a waveform that is at one level for a short
time and at another level for the rest of the period, characterized by
its \textbf{duty cycle} (fraction of the period spent ``high''). Digital
data and radar are naturally pulse-like.
\end{itemize}

\section{Fourier's Theorem}
\label{sec:fourier}

The mathematician Joseph Fourier proved a result of enormous consequence
for radio engineering: \emph{any} periodic waveform, no matter its shape,
can be expressed as a sum (in general, an infinite sum) of sine and
cosine waves at the fundamental frequency and its integer multiples
(\textbf{harmonics}):
\[
x(t) = a_0 + \sum_{n=1}^{\infty} \left[a_n \cos(n\omega_0 t) + b_n \sin(n\omega_0 t)\right]
\]
where $\omega_0 = 2\pi f_0$ is the fundamental angular frequency and
$a_0, a_n, b_n$ are coefficients determined by the waveform's shape. This
means a waveform's \emph{shape} and its \emph{frequency content} are two
views of the same underlying signal, related by the \textbf{Fourier
transform} (for non-periodic or more general signals) or the
\textbf{Fourier series} (for periodic signals, as written above).

\begin{keyidea}
Every non-sinusoidal periodic signal occupies more than one frequency: a
fundamental plus a ladder of harmonics. A signal's bandwidth --- how much
of the radio spectrum it actually occupies --- is fundamentally a
statement about its Fourier content, not just its nominal ``frequency.''
\end{keyidea}

\subsection{Square Wave Harmonics}

A symmetric square wave of fundamental frequency $f_0$ and peak amplitude
$A$ contains \emph{only odd} harmonics, with amplitude falling off as
$1/n$:
\[
x(t) = \frac{4A}{\pi}\left[\sin(\omega_0 t) + \frac{1}{3}\sin(3\omega_0 t) + \frac{1}{5}\sin(5\omega_0 t) + \cdots\right]
\]
This is why a keyed CW signal or a switching power supply, even though it
is nominally operating at one frequency, actually radiates energy at that
frequency's odd multiples too --- a real-world source of harmonic
interference (Section~\ref{sec:harmonic-interference}) unless filtered.

\begin{examplebox}[Predicting a harmonic frequency]
A square-wave clock oscillates at 7.023\,MHz. At what frequency would you
expect to find its strongest spurious harmonic, and in which amateur band
might it fall?

\emph{Solution.} The strongest harmonic of a square wave (after the
fundamental) is the third harmonic: $3 \times 7.023 = 21.069\,\text{MHz}$,
which falls inside the 15\,m amateur band (21.000--21.450\,MHz in most
IARU Region 1/2/3 allocations) --- a classic real-world case for
low-pass filtering a keyed oscillator or transmitter output.
\end{examplebox}

\section{Bandwidth and Rise Time}

A perfectly square edge (instantaneous transition) requires
\emph{infinite} bandwidth to reproduce exactly, because it needs
harmonics out to $n \to \infty$ to build the sharp corner. Real circuits
have finite bandwidth, so real square waves always have a finite
\textbf{rise time} $t_r$ (time to transition, conventionally measured
from 10\% to 90\% of full amplitude). A useful rule of thumb links rise
time and the bandwidth $B$ needed to reproduce it reasonably faithfully:
\[
B \approx \frac{0.35}{t_r}.
\]
This relationship explains why fast-keyed digital or CW signals, or
poorly shaped keying waveforms, splatter energy into adjacent parts of
the band --- an important practical and courtesy consideration when
operating, and the reason well-designed transmitters deliberately shape
(``soften'') their keying edges.

\begin{safetynote}[Key clicks and splatter]
Overly fast keying transitions (``hard keying'') generate broadband
clicks that can interfere with stations operating hundreds of kilohertz
away. Modern transmitters shape the CW envelope with a short rise/fall
time (typically several milliseconds) specifically to keep the emitted
bandwidth narrow. This is both good operating practice and, in most
countries, a regulatory requirement on out-of-band emissions.
\label{sec:harmonic-interference}
\end{safetynote}

\section{DC Offset and the Mean Value}

The $a_0$ term in the Fourier series represents the waveform's
\textbf{DC component} --- its average value over a full cycle. A sine
wave that swings symmetrically about zero has $a_0 = 0$; a keyed pulse
train that is ``on'' more than it is ``off'' has a nonzero DC component
equal to (peak value) $\times$ (duty cycle). This matters practically
when a waveform passes through a transformer or capacitor-coupled stage,
neither of which passes a DC component --- the shape of the waveform will
change (its average will be forced back to zero) even though its AC
content is otherwise preserved.

\section{Spectral (Frequency-Domain) View}

Plotting a signal's amplitude at each frequency, rather than its
instantaneous value over time, gives its \textbf{spectrum} --- exactly
what a spectrum analyzer or waterfall display shows. A pure sine wave
appears as a single vertical line at its frequency; a square wave appears
as a line at the fundamental plus progressively smaller lines at each odd
harmonic; a burst of noise appears as energy spread continuously across a
range of frequencies. Learning to think in both the \textbf{time domain}
(a scope trace) and the \textbf{frequency domain} (a spectrum display)
simultaneously is one of the most valuable skills a radio amateur can
develop, and is essential for understanding modulation
(Chapter~\ref{ch:modulation}).

\begin{quickreview}
\item Any periodic waveform is a sum of sine/cosine harmonics of its
fundamental frequency (Fourier series).
\item A symmetric square wave contains only odd harmonics, falling off as
$1/n$ in amplitude.
\item Sharper edges (shorter rise time) require more bandwidth,
approximately $B \approx 0.35/t_r$.
\item A waveform's DC component is its Fourier $a_0$ term --- the
average value over one cycle.
\item Time-domain and frequency-domain views describe the same signal
from two complementary perspectives.
\end{quickreview}

\begin{practiceproblems}
\item List the first three nonzero harmonics (as multiples of $f_0$) of a
symmetric square wave.
\item A 3.573\,MHz CW oscillator's third harmonic falls near which
amateur band?
\item A digital keying circuit has a rise time of 1\,ms. Estimate the
bandwidth this implies, and comment on whether this is adequate for
narrow, splatter-free CW.
\item Explain why a capacitor-coupled audio stage will not faithfully
pass a waveform with a strong DC offset.
\item Sketch (in words) how the spectrum of a triangle wave would differ
from that of a square wave of the same fundamental frequency.
\end{practiceproblems}
